\subsection{Regresión Ridge: un paso}

En \texttt{ridge.py} se añadió el rezago de una observación de la generación, \texttt{Gen\_lag1}, a los cinco predictores de calendario. Se eliminaron las filas con valores faltantes mediante \texttt{dropna} y se tomó el primer 70\,\% del conjunto resultante. Cada ventana contiene 13 registros: 12 para ajustar y uno para evaluar. También se reajusta el modelo en cada desplazamiento.

Se utilizó \texttt{RidgeCV} con los candidatos $\alpha\in\{0.001,0.01,0.1,1,10\}$ y tres particiones de validación interna. La penalización corresponde a
\begin{equation}
(\hat{\beta}_0,\hat{\boldsymbol{\beta}})
=\arg\min_{\beta_0,\boldsymbol{\beta}}
\left\{
\sum_{i=1}^{12}
\left(y_i-\beta_0-\boldsymbol{x}_i^{\mathsf T}\boldsymbol{\beta}\right)^2
+\alpha\|\boldsymbol{\beta}\|_2^2
\right\}.
\end{equation}
Este programa no aplica estandarización a las entradas. Las tres particiones internas tampoco constituyen una validación de origen móvil: no se especifica un particionador temporal. La salida comprende MAE, RMSE, $R^2$ y \texttt{prediccion\_vs\_real\_ridge.png}. La comparación con OLS involucra simultáneamente cambios de predictores, tamaño de entrenamiento y escalamiento; por ello, no permite aislar el efecto de la penalización.
